% ============================================================
%  GLOSARIO Y ABREVIATURAS
%  ------------------------------------------------------------
%  Va al final de los preliminares, después del índice de
%  figuras (orden exigido por la guía §2.2), y se incluye solo
%  si el interruptor "mostrarOpcionales" está encendido.
%  Se incorpora únicamente cuando el proyecto usa términos
%  técnicos, siglas o conceptos que requieren aclaración: si el
%  proyecto no los usa, basta con borrar el cuerpo del entorno
%  y el capítulo desaparece del documento sin dejar hojas en
%  blanco (no hay entradas vacías porque la lista está completa).
% ============================================================

\chapter*{GLOSARIO Y ABREVIATURAS}
\addcontentsline{toc}{chapter}{GLOSARIO Y ABREVIATURAS}
\begingroup           % \interlineadoResumen no debe filtrarse al resto del documento
\interlineadoResumen
\begin{description}[leftmargin=3cm,style=nextline,font=\bfseries]
  \item[Autocorrelación espacial] Tendencia de las unidades cercanas en el espacio a
    presentar valores similares de una variable.
  \item[B0, B1, B2, B3] Líneas base de tasa, ordenadas de menor a mayor complejidad: global
    (B0), por municipio (B1), por anillo de distancia (B2) y por vecindad H3 (B3). La letra B
    viene de \textit{baseline}.
  \item[Bloque espacial] Celda H3 de resolución 6 que agrupa unas 49 zonas de resolución 8;
    unidad de partición para la validación y la prueba.
  \item[Brecha] Diferencia, en una escala comparable entre zonas, entre las consultas
    observadas y las esperadas; se mide con el residuo de Pearson,
    $(\text{observado} - \text{esperado}) / \sqrt{\text{esperado}}$.
  \item[Calibración] Cociente entre la suma de las predicciones y la suma de los valores
    observados.
  \item[CRISP-DM] \textit{Cross-Industry Standard Process for Data Mining}; metodología de
    seis fases para proyectos de Ciencia de Datos.
  \item[D\textsuperscript{2}] Fracción de la devianza explicada respecto de predecir el
    promedio del conjunto evaluado.
  \item[Devianza de Poisson] Medida de error para datos de conteo que evalúa el error en
    términos relativos.
  \item[$\delta$ de Cliff] Medida no paramétrica del tamaño de la diferencia
    entre dos grupos, en escala de $-1$ a $+1$.
  \item[EDA] Análisis Exploratorio de Datos (\textit{Exploratory Data Analysis}).
  \item[Exposición] Magnitud a la que se espera que el conteo sea proporcional; en este
    trabajo, la población de la zona.
  \item[Fuga de información] Uso inadvertido, durante el entrenamiento, de información de
    la respuesta o del conjunto de prueba (\textit{data leakage}).
  \item[GLM] Modelo lineal generalizado (\textit{Generalized Linear Model}).
  \item[GTFS] \textit{General Transit Feed Specification}; estándar abierto para describir
    rutas, paradas y horarios de transporte público.
  \item[H3] Sistema de indexación geoespacial hexagonal y jerárquico.
  \item[I de Moran] Estadístico de autocorrelación espacial global.
  \item[Índice de Jaccard] Proporción de elementos compartidos entre dos conjuntos; aquí,
    entre dos listas de zonas prioritarias.
  \item[Kontur Population] Conjunto de datos abiertos de población estimada sobre
    hexágonos H3.
  \item[Línea base] Punto de referencia inicial: una estimación deliberadamente simple que
    cualquier técnica más elaborada debe superar para justificar su complejidad.
  \item[LISA] Indicadores locales de asociación espacial (\textit{Local Indicators of
    Spatial Association}).
  \item[M1, M2, M3] Modelos que estiman parámetros a partir de predictores: GLM de Poisson
    (M1), GLM binomial negativa (M2) y \textit{gradient boosting} con pérdida de Poisson
    (M3). La letra M viene de \textit{modelo}.
  \item[MAE] Error absoluto medio.
  \item[MR] Mejora relativa frente al modelo de referencia (tasa global, B0):
    $1 - D_{\mathrm{modelo}} / D_{\mathrm{referencia}}$, fracción de la devianza de B0 que
    reduce una técnica.
  \item[Offset] Término del modelo con coeficiente fijado en 1; en este trabajo,
    $\log(\text{población})$.
  \item[Paratránsito] Transporte público operado por sindicatos y cooperativas con rutas
    flexibles y sin horarios fijos.
  \item[Regla de un error estándar] Criterio de selección: se elige la técnica más simple
    cuya devianza media no supera la mejor devianza media más un error estándar de esta.
  \item[Sobredispersión] Condición en la que la varianza de un conteo supera a su media.
  \item[Spearman, correlación de] Correlación basada en rangos entre dos variables u
    ordenaciones.
  \item[Zona] Celda H3 de resolución 8 (unos 0,74~km\textsuperscript{2}); unidad de análisis
    del proyecto.
\end{description}
\endgroup
\clearpage
